\documentclass[10pt,a4paper]{amsart}
\usepackage[margin=19mm]{geometry}
\usepackage{amsmath,amssymb,mathtools}
\usepackage[scr=rsfs,cal=euler]{mathalfa}
\usepackage{xfrac}
\usepackage[unicode,hidelinks,bookmarks=false]{hyperref}
\newcommand{\ie}{i.e.\ }
\newcommand{\cf}{cf.\ }
\newcommand{\resp}{respectively\ }
\newcommand{\Subsection}{\subsection}
\providecommand{\nobreakdash}{\nobreak\hbox{-}}
\setlength{\emergencystretch}{2em}
\multlinegap0pt
\usepackage{amssymb}
\usepackage{mathtools}
\usepackage[shortlabels]{enumitem}
\newtheorem{lemma}{Lemma}
\newcommand{\norm}[1]{\left\lVert #1\right\rVert}
\title{Finite-Window Recursive Audit Chains for Navier--Stokes Generated Packages}
\author{Runlong Yu}
\date{Source: arXiv:2606.20899v1 (CC BY 4.0)}
\begin{document}
\maketitle

\noindent\emph{Source and licence.} Finite-Window Recursive Audit Chains for Navier--Stokes Generated Packages. Version arXiv:2606.20899v1 (CC BY 4.0), distributed by arXiv under the Creative Commons Attribution 4.0 International licence, http://creativecommons.org/licenses/by/4.0/, as recorded in the arXiv licence metadata for this exact version and shown on the abstract page https://arxiv.org/abs/2606.20899v1. Full licence notice: \texttt{licenses/arxiv-2606.20899-CC-BY-4.0.txt}.\par\smallskip

\noindent\textit{Editorial scope.} Counted unit: the article's lemma lem:ckn-closeness-smallness (source0.tex lines 4266--4304). The article's complete original proof of it is delivered with it (source0.tex lines 4306--4391, one \texttt{proof} environment of 86 lines). No mathematics is written, repaired or re-derived: the statement and the proof are the article's own lines, in the article's own notation, and no retained line is substituted or re-flowed.  No further source-local context is needed: every cross-reference of every retained line resolves inside the two delivered spans.  Nothing was omitted from any retained span, so the package carries no omission record.  No retained line contains a citation, so the wrapper carries no bibliography and no bibliography entry.  No retained line contains a figure environment, an image-inclusion command or drawing code, so no image is delivered and the package declares no asset dependency.  Editorial formatting only, in addition: an AMS \texttt{amsart} preamble that restates the article's own declarations and macros used by the retained lines, namely the environments \texttt{lemma} on separate counters, together with the standard packages those lines need.  The retained environments print in this document as Lemma 1 in order of appearance, whereas the article prints the same items with the numbering of its own full document.  The complete native source of the article as distributed in the arXiv e-print for this exact version is bundled unchanged as \texttt{sources/203245/source0.tex}.
\begin{lemma}[CKN smallness from CKN-compatible quotient closeness]
\label{lem:ckn-closeness-smallness}
Assume there exists \(\zeta\in\Gamma_{\Lambda,k}^{\mathrm{CKN}}\) such that
\[
  \norm{u_k-u_\zeta}_{L^3(Q_\vartheta)}
  +
  \inf_{a=a(s)}
  \norm{p_k-p_\zeta-a(s)}_{L^{3/2}(Q_\vartheta)}
  \le
  \tau.
\]
Then there is a constant \(C_\vartheta\) such that
\[
  \Phi_k(\vartheta)
  \le
  2\Phi_\zeta(\vartheta)
  +
  C_\vartheta(\tau^3+\tau^{3/2}).
\]
In particular, since
\(\Phi_\zeta(\vartheta)\le\varepsilon_{\mathrm{CKN}}/4\), one also has the
coarser bound
\[
  \Phi_k(\vartheta)
  \le
  C_\vartheta
  \left[
  \frac{\varepsilon_{\mathrm{CKN}}}{4}
  +
  \tau^3+\tau^{3/2}
  \right].
\]
Consequently there exists \(\tau_{\mathrm{CKN}}>0\) such that
\[
  \tau\le\tau_{\mathrm{CKN}}
  \quad\Longrightarrow\quad
  \Phi_k(\vartheta)<\varepsilon_{\mathrm{CKN}}.
\]
\end{lemma}

\begin{proof}
Fix \(\eta_0>0\).  Since the pressure distance contains an infimum over
time-dependent constants, choose \(a(s)\) such that
\[
  \norm{u_k-u_\zeta}_{L^3(Q_\vartheta)}
  +
  \norm{p_k-p_\zeta-a(s)}_{L^{3/2}(Q_\vartheta)}
  \le
  \tau+\eta_0.
\]
For every \(r>1\) and every \(\alpha>0\),
\[
  |A+B|^r\le (1+\alpha)|A|^r+C_{r,\alpha}|B|^r .
\]
Using this with \(r=3\), \(A=u_\zeta\), and
\(B=u_k-u_\zeta\), and choosing \(\alpha\le1\), gives
\[
  \int_{Q_\vartheta}|u_k|^3
  \le
  2\int_{Q_\vartheta}|u_\zeta|^3
  +
  C\norm{u_k-u_\zeta}_{L^3(Q_\vartheta)}^3.
\]

For pressure, write
\[
  q:=p_k-p_\zeta-a(s).
\]
Since spatial mean subtraction removes the time-dependent pressure constant,
\[
  p_k-(p_k)_{B_\vartheta}(s)
  =
  p_\zeta-(p_\zeta)_{B_\vartheta}(s)
  +
  q-(q)_{B_\vartheta}(s).
\]
Applying the same perturbative inequality with \(r=3/2\) yields
\[
  \int_{Q_\vartheta}|p_k-(p_k)_{B_\vartheta}(s)|^{3/2}
  \le
  2\int_{Q_\vartheta}
  |p_\zeta-(p_\zeta)_{B_\vartheta}(s)|^{3/2}
  +
  C\int_{Q_\vartheta}|q-(q)_{B_\vartheta}(s)|^{3/2}.
\]
The mean-subtraction map is bounded on \(L^{3/2}(B_\vartheta)\), uniformly for
the fixed ball \(B_\vartheta\), so
\[
  \int_{Q_\vartheta}|q-(q)_{B_\vartheta}(s)|^{3/2}
  \le
  C_\vartheta\norm{q}_{L^{3/2}(Q_\vartheta)}^{3/2}.
\]
Combining the velocity and pressure estimates and using the choice of
\(a(s)\) gives
\[
  \Phi_k(\vartheta)
  \le
  2\Phi_\zeta(\vartheta)
  +
  C_\vartheta\bigl((\tau+\eta_0)^3+(\tau+\eta_0)^{3/2}\bigr).
\]
Letting \(\eta_0\downarrow0\) proves the first estimate.  The second estimate
follows from \(\Phi_\zeta(\vartheta)\le\varepsilon_{\mathrm{CKN}}/4\).

Choose \(\tau_{\mathrm{CKN}}>0\) so small that
\[
  C_\vartheta
  \left(
  \tau_{\mathrm{CKN}}^3+\tau_{\mathrm{CKN}}^{3/2}
  \right)
  <
  \frac{\varepsilon_{\mathrm{CKN}}}{2}.
\]
Then the first estimate gives
\[
  \Phi_k(\vartheta)
  <
  \frac{\varepsilon_{\mathrm{CKN}}}{2}
  +
  \frac{\varepsilon_{\mathrm{CKN}}}{2}
  =
  \varepsilon_{\mathrm{CKN}},
  \qquad
  \tau\le\tau_{\mathrm{CKN}}.
\]
\end{proof}

\end{document}
